% ═══════════════════════════════════════════════════════════════════════
\section{Overleaf : \texorpdfstring{\LaTeX}{LaTeX} dans le navigateur}
% ═══════════════════════════════════════════════════════════════════════

\begin{frame}{Démarrer sur Overleaf}
  \small
  \begin{enumerate}
    \item Aller sur \texttt{overleaf.com} et créer un compte gratuit.
    \item Nouveau projet (New Project) $\to$ Projet vide (Blank Project),
      lui donner un nom.
    \item Pour partir d'un fichier fourni : Nouveau projet $\to$
      Téléverser un projet (Upload Project), ou téléverser le
      \texttt{.tex} dans un projet existant.
    \item Rien à installer : la compilation se fait sur les serveurs
      d'Overleaf.
  \end{enumerate}
  \vspace{4pt}
  \begin{Rem}{Bon à savoir}
    \small Le compilateur par défaut est pdfLaTeX : c'est celui que
    supposent nos fichiers d'exemple. On peut le changer dans le
    menu (Menu $\to$ Compilateur), mais ce n'est pas nécessaire
    aujourd'hui.
  \end{Rem}
\end{frame}

\begin{frame}{L'interface}
  \FigBox[0.6\textheight]{%
    \begin{tikzpicture}[x=1cm, y=1cm, font=\scriptsize\sffamily]
      \draw[sectioncolor, thick, rounded corners=2pt] (0,0) rectangle (13,6.2);
      \fill[sectioncolor!12] (0,5.6) rectangle (13,6.2);
      \node[anchor=west] at (0.2,5.9) {Menu \quad Mon projet};
      \node[fill=vertpiece!80!black, text=white, rounded corners=2pt,
        inner sep=3pt] at (8.3,5.9) {Recompiler};
      \draw[sectioncolor!60] (2.6,0) -- (2.6,5.6);
      \draw[sectioncolor!60] (7.4,0) -- (7.4,5.6);
      \node[anchor=north west, align=left] at (0.15,5.45)
        {Fichiers\\[2pt] \texttt{main.tex}\\ \texttt{figure.png}};
      \node[anchor=north west, align=left, font=\scriptsize\ttfamily] at (2.75,5.45)
        {\textbackslash documentclass\{article\}\\
         \textbackslash usepackage[french]\{babel\}\\[2pt]
         \textbackslash begin\{document\}\\
         Bonjour !\\
         \textbackslash end\{document\}};
      \fill[white] (7.7,0.3) rectangle (12.7,5.3);
      \draw[gray!60] (7.7,0.3) rectangle (12.7,5.3);
      \node at (10.2,4.6) {Bonjour !};
      \node[firebrick3, align=center] at (1.3,-0.5) {1. les fichiers\\ du projet};
      \node[firebrick3, align=center] at (5,-0.5) {2. l'éditeur :\\ le code \texttt{.tex}};
      \node[firebrick3, align=center] at (10.2,-0.5) {3. l'aperçu :\\ le PDF compilé};
      \node[firebrick3, align=center] at (10.4,6.65) {4. compiler (Ctrl+Entrée,\\ Cmd+Entrée sur Mac)};
    \end{tikzpicture}}
  \small On modifie le code à gauche, on recompile, on regarde le PDF à
  droite. Les erreurs s'affichent près du bouton Recompiler.
\end{frame}

\begin{frame}{Le cycle de travail}
  \FigBox[0.32\textheight]{%
    \begin{tikzpicture}[
        etape/.style={draw=sectioncolor, fill=thmboxbg, rounded corners=2pt,
          thick, align=center, font=\small, minimum height=1.1cm,
          text width=2.8cm}]
      \node[etape] (a) at (0,0) {Écrire ou modifier\\le code \texttt{.tex}};
      \node[etape] (b) at (4.2,0) {Compiler\\(pdfLaTeX)};
      \node[etape] (c) at (8.4,0) {Lire le PDF\\(ou l'erreur)};
      \draw[fleche] (a) -- (b);
      \draw[fleche] (b) -- (c);
      \draw[fleche] (c.south) to[out=-150, in=-30] (a.south);
    \end{tikzpicture}}
  \small
  \begin{itemize}
    \item Compiler souvent : après chaque petite modification.
    \item Une erreur à la fois : la première erreur signalée est presque
      toujours la vraie ; les suivantes en découlent souvent.
    \item Le PDF reste celui de la dernière compilation réussie : vérifier
      qu'il a bien changé.
  \end{itemize}
\end{frame}

\begin{frame}{À vous}
  \begin{Exercice}{Un premier projet}
    \small
    \begin{enumerate}
      \item Créez un projet sur Overleaf et téléversez le fichier
        \texttt{01-minimal.tex}.
      \item Compilez-le : vous devez obtenir une page avec trois sections.
      \item Remplacez « Prénom Nom » par votre nom, recompilez.
      \item Ajoutez un quatrième élément à la liste, recompilez.
    \end{enumerate}
  \end{Exercice}
\end{frame}
